%
% 11-problem.tex -- Problemstellung
%
% (c) 2026 Prof Dr Andreas Müller
%
\subsection{Problemstellung}
Die von Bernoulli gelöste Aufgabe ist, die Stammfunktion einer
rationalen Funktion zu finden, für die ausserdem angenommen wird,
dass eine Faktorisierung des Nenners in Linearfaktoren bekannt ist.
Der Fundamentalsatz der Algebra garantiert, dass eine solche
Faktorisierung immer exisiert, wenngleich er keine Methode
bereitstellt, die Faktorisierung zu finden.

\begin{aufgabe}
Gegeben ist eine rationale Funktion $f(z) \in \mathbb{C}(z)$ mit
komplexen Koeffizienten.
\index{rationale Funktion}%
Finde eine Stammfunktion $F$.
\index{Stammfunktion}
\end{aufgabe}

Die rationale Funktion wird als gekürzter Bruch
\[
f(z)
=
\frac{p(z)}{q(z)}
\qquad\text{und}\qquad
p(z),q(z)\in \mathbb{Q}(z)
\]
geschrieben, wobei $p(z)$ und $q(z)$ teilerfremd sind.
Es wird angenommen, dass der Nenner in ein Produkt
\[
q(z)
=
(z-\alpha_1)\cdot\ldots\cdot (z-\alpha_n)
=
\prod_{k=1}^n (z-\alpha_k)
\]
faktorisiert werden kann.

Bernoullis Strategie ist, für die rationale Funktion erst
ein Partialbruchzerlegung zu finden, was im
Abschnitt~\ref{buch:ratint:bernoulli:subsection:partialbruch}
durchgeführt werden soll.
Er schreibt den Integranden als Summe
\begin{align*}
f(z)
=
s(z)
+
\sum_{k=1}^n\sum_{l=1}^{n_k}\frac{A_k^l}{(z-\alpha_k)^l},
\end{align*}
wobei $s(z)$ ein Polynom und $\alpha_k$ die Nullstellen
des Nenners von $f(z)$ sind. 
Um eine Stammfunktion zu finden, müssen Stammfunktionen der
einzelnen Summanden gefunden werden.
Der Polynomterm ist dabei das kleinste Problem, da die einzelnen
Summanden $z^n$ die Stammfunktion $z^{n+1}/(n+1)$ haben.
Etwas mehr Arbeit geben die Integrale der Brüche, die aber ebenfalls
wohlbekannt sind und in
Gleichung~\eqref{buch:ratint:bernoulli:partial:eqn:stammfunktionen}
angegeben werden.

